% GENERATED FILE. Edit canonical lessons, then run npm run generate:books.
% canonical-source-hash: fnv1a64:b238659780535c12
% canonical-lessons: SA-C350-bandhanam, SA-C350-muktih, SA-C350-nirdharayati, SA-C350-pratyarpayati, SA-C350-prayascittam

\chapter{Release}
\label{ch:sa-release}
\hlchaptermodality{350}

\begin{chapteropening}
\textbf{By the end of this chapter:} \emph{I can say who was set free and what made amends.}

Complete the last lesson of chapter 350: I can say who was set free and what made amends.
\end{chapteropening}

\section[\texorpdfstring{\sk{बन्धनम्} (bandhanam)}{बन्धनम् (bandhanam)}]{\sk{बन्धनम्} (bandhanam) — captivity; a bond}

\label{lesson:SA-C350-bandhanam}



\begin{quote}
Before the new one: say the Sanskrit for restrains, then the Sanskrit for an oath.
\end{quote}

\subsection*{You'll want to know: \sk{बन्धनम्}}
\textbf{\sk{बन्धनम्}} — \emph{bandhanam} — \textquotedblleft{}captivity; a bond\textquotedblright{}.

\begin{quote}
\sk{पक्षी} \sk{बन्धनं} \sk{न} \sk{इच्छति।}
\end{quote}

\begin{quote}
\emph{pakṣī bandhanaṃ na icchati.}
\end{quote}

\begin{quote}
A bird does not want captivity.
\end{quote}

A neuter noun.

\begin{cousinweb}
\textbf{\sk{बन्धनम्}} (\emph{bandhanam}) comes from the root √bandh, \textquotedblleft{}to bind\textquotedblright{}; MW gives \emph{bandhana} \textquotedblleft{}the act of binding, tying, fettering\textquotedblright{} (MW).
\end{cousinweb}

\begin{grammarlens}[title={What you've built}]
One more word for giving reasons and opinions.
\end{grammarlens}

\subsection*{Your turn}
\begin{itemize}
\raggedright
  \item \emph{Say it:} \emph{bandhanam}
  \item \emph{Say it:} \emph{bandhanam}, once more
  \item \emph{Say it:} \emph{pakṣī bandhanaṃ na icchati.} — A bird does not want captivity.
  \item \emph{Recall:} say the Sanskrit for restrains, then the Sanskrit for an oath, then say \emph{bandhanam} again
\end{itemize}

\begin{tcolorbox}[breakable,colback=teal!4,colframe=teal!35!black,title={Before you move on}]
What does \sk{बन्धनम्} mean? (\textquotedblleft{}Captivity; a bond\textquotedblright{}.) Say it once more.
\end{tcolorbox}
\section[\texorpdfstring{\sk{मुक्तिः} (muktiḥ)}{मुक्तिः (muktiḥ)}]{\sk{मुक्तिः} (muktiḥ) — release; freedom}

\label{lesson:SA-C350-muktih}



\begin{quote}
Before the new one: say the Sanskrit for an oath, then the Sanskrit for captivity.
\end{quote}

\subsection*{You'll want to know: \sk{मुक्तिः}}
\textbf{\sk{मुक्तिः}} — \emph{muktiḥ} — \textquotedblleft{}release; freedom\textquotedblright{}.

\begin{quote}
\sk{बन्दी} \sk{मुक्तिं} \sk{याचते।}
\end{quote}

\begin{quote}
\emph{bandī muktiṃ yācate.}
\end{quote}

\begin{quote}
The captive begs for release.
\end{quote}

A feminine noun.

\begin{cousinweb}
\textbf{\sk{मुक्तिः}} (\emph{muktiḥ}) comes from the root √muc, \textquotedblleft{}to set free\textquotedblright{}; MW gives \emph{mukti} \textquotedblleft{}setting or becoming free, release, liberation\textquotedblright{} (MW).
\end{cousinweb}

\begin{grammarlens}[title={What you've built}]
One more word for giving reasons and opinions.
\end{grammarlens}

\subsection*{Your turn}
\begin{itemize}
\raggedright
  \item \emph{Say it:} \emph{muktiḥ}
  \item \emph{Say it:} \emph{muktiḥ}, once more
  \item \emph{Say it:} \emph{bandī muktiṃ yācate.} — The captive begs for release.
  \item \emph{Recall:} say the Sanskrit for an oath, then the Sanskrit for captivity, then say \emph{muktiḥ} again
\end{itemize}

\begin{tcolorbox}[breakable,colback=teal!4,colframe=teal!35!black,title={Before you move on}]
What does \sk{मुक्तिः} mean? (\textquotedblleft{}Release; freedom\textquotedblright{}.) Say it once more.
\end{tcolorbox}
\section[\texorpdfstring{\sk{निर्धारयति} (nirdhārayati)}{निर्धारयति (nirdhārayati)}]{\sk{निर्धारयति} (nirdhārayati) — settles (a question); determines}

\label{lesson:SA-C350-nirdharayati}



\begin{quote}
Before the new one: say the Sanskrit for captivity, then the Sanskrit for release.
\end{quote}

\subsection*{You'll want to know: \sk{निर्धारयति}}
\textbf{\sk{निर्धारयति}} — \emph{nirdhārayati} — \textquotedblleft{}settles (a question); determines\textquotedblright{}.

\begin{quote}
\sk{न्यायाधीशः} \sk{दण्डं} \sk{निर्धारयति।}
\end{quote}

\begin{quote}
\emph{nyāyādhīśaḥ daṇḍaṃ nirdhārayati.}
\end{quote}

\begin{quote}
The judge settles the punishment.
\end{quote}

Past (he …): \sk{सः} \sk{दण्डं} \sk{निर्धारितवान्}, \emph{saḥ daṇḍaṃ nirdhāritavān}.

\begin{cousinweb}
\textbf{\sk{निर्धारयति}} (\emph{nirdhārayati}) is built on \emph{nir}, \textquotedblleft{}out\textquotedblright{}, and the root √dhṛ, \textquotedblleft{}to hold\textquotedblright{}, in its causative form; MW gives \emph{nir-dhārita} \textquotedblleft{}determined, ascertained, settled\textquotedblright{} (MW).
\end{cousinweb}

\begin{grammarlens}[title={What you've built}]
One more word for giving reasons and opinions.
\end{grammarlens}

\subsection*{Your turn}
\begin{itemize}
\raggedright
  \item \emph{Say it:} \emph{nirdhārayati}
  \item \emph{Say it:} \emph{nirdhārayati}, once more
  \item \emph{Say it:} \emph{nyāyādhīśaḥ daṇḍaṃ nirdhārayati.} — The judge settles the punishment.
  \item \emph{Recall:} say the Sanskrit for captivity, then the Sanskrit for release, then say \emph{nirdhārayati} again
\end{itemize}

\begin{tcolorbox}[breakable,colback=teal!4,colframe=teal!35!black,title={Before you move on}]
What does \sk{निर्धारयति} mean? (\textquotedblleft{}Settles (a question); determines\textquotedblright{}.) Say it once more.
\end{tcolorbox}
\section[\texorpdfstring{\sk{प्रत्यर्पयति} (pratyarpayati)}{प्रत्यर्पयति (pratyarpayati)}]{\sk{प्रत्यर्पयति} (pratyarpayati) — gives back; restores}

\label{lesson:SA-C350-pratyarpayati}



\begin{quote}
Before the new one: say the Sanskrit for release, then the Sanskrit for settles (a question).
\end{quote}

\subsection*{You'll want to know: \sk{प्रत्यर्पयति}}
\textbf{\sk{प्रत्यर्पयति}} — \emph{pratyarpayati} — \textquotedblleft{}gives back; restores\textquotedblright{}.

\begin{quote}
\sk{चौरः} \sk{धनं} \sk{प्रत्यर्पयति।}
\end{quote}

\begin{quote}
\emph{cauraḥ dhanaṃ pratyarpayati.}
\end{quote}

\begin{quote}
The thief gives back the money.
\end{quote}

Past (he …): \sk{सः} \sk{स्यूतं} \sk{प्रत्यर्पितवान्}, \emph{saḥ syūtaṃ pratyarpitavān}.

\begin{cousinweb}
\textbf{\sk{प्रत्यर्पयति}} (\emph{pratyarpayati}) is built on \emph{prati}, \textquotedblleft{}back\textquotedblright{}, and the root √ṛ, \textquotedblleft{}to go\textquotedblright{}, in its causative form \emph{arpaya}, \textquotedblleft{}to hand over\textquotedblright{}; MW gives \emph{praty-arpita} \textquotedblleft{}restored\textquotedblright{} (MW).
\end{cousinweb}

\begin{grammarlens}[title={What you've built}]
One more word for giving reasons and opinions.
\end{grammarlens}

\subsection*{Your turn}
\begin{itemize}
\raggedright
  \item \emph{Say it:} \emph{pratyarpayati}
  \item \emph{Say it:} \emph{pratyarpayati}, once more
  \item \emph{Say it:} \emph{cauraḥ dhanaṃ pratyarpayati.} — The thief gives back the money.
  \item \emph{Recall:} say the Sanskrit for release, then the Sanskrit for settles (a question), then say \emph{pratyarpayati} again
\end{itemize}

\begin{tcolorbox}[breakable,colback=teal!4,colframe=teal!35!black,title={Before you move on}]
What does \sk{प्रत्यर्पयति} mean? (\textquotedblleft{}Gives back; restores\textquotedblright{}.) Say it once more.
\end{tcolorbox}
\section[\texorpdfstring{\sk{प्रायश्चित्तम्} (prāyaścittam)}{प्रायश्चित्तम् (prāyaścittam)}]{\sk{प्रायश्चित्तम्} (prāyaścittam) — atonement; amends}

\label{lesson:SA-C350-prayascittam}



\begin{quote}
Before the new one: say the Sanskrit for settles (a question), then the Sanskrit for gives back.
\end{quote}

\subsection*{You'll want to know: \sk{प्रायश्चित्तम्}}
\textbf{\sk{प्रायश्चित्तम्}} — \emph{prāyaścittam} — \textquotedblleft{}atonement; amends\textquotedblright{}.

\begin{quote}
\sk{दोषस्य} \sk{प्रायश्चित्तं} \sk{करोतु।}
\end{quote}

\begin{quote}
\emph{doṣasya prāyaścittaṃ karotu.}
\end{quote}

\begin{quote}
Make amends for the fault.
\end{quote}

A neuter noun.

\begin{cousinweb}
\textbf{\sk{प्रायश्चित्तम्}} (\emph{prāyaścittam}) is built from \emph{prāyas}, \textquotedblleft{}for the most part\textquotedblright{}, and \emph{citta}, \textquotedblleft{}thought\textquotedblright{} (MW records the explanations \textquotedblleft{}predominant thought\textquotedblright{} and \textquotedblleft{}thought of death\textquotedblright{}); MW gives \emph{prāyaś-citta} \textquotedblleft{}atonement, expiation, amends\textquotedblright{} (MW).
\end{cousinweb}

\begin{grammarlens}[title={What you've built}]
One more word for giving reasons and opinions.
\end{grammarlens}

\subsection*{Your turn}
\begin{itemize}
\raggedright
  \item \emph{Say it:} \emph{prāyaścittam}
  \item \emph{Say it:} \emph{prāyaścittam}, once more
  \item \emph{Say it:} \emph{doṣasya prāyaścittaṃ karotu.} — Make amends for the fault.
  \item \emph{Recall:} say the Sanskrit for settles (a question), then the Sanskrit for gives back, then say \emph{prāyaścittam} again
\end{itemize}

\begin{tcolorbox}[breakable,colback=teal!4,colframe=teal!35!black,title={Before you move on}]
What does \sk{प्रायश्चित्तम्} mean? (\textquotedblleft{}Atonement; amends\textquotedblright{}.) Say it once more.
\end{tcolorbox}
